This project aims to study, design, and develop an automated Category-Specific News Classification and NLP Analytics Pipeline. The system covers an end-to-end workflow including multi-source news aggregation, content extraction, summarization, category classification, trend analytics, and real-time delivery. The architecture is engineered according to SOLID and GRASP software design principles with clearly decoupled layers. The backend is implemented in Python using FastAPI, supporting automated extraction from both standard RSS feeds and dynamic JavaScript-heavy websites via a Playwright headless browser instance. A multi-tier caching strategy (combining in-memory LRU caching and HTTP 304 conditional cache) is employed to minimize network overhead and redundant processing. For content summarization and sentiment analysis, the system integrates with Large Language Models (LLMs) via the KKU Gen.AI API to generate concise summaries, bullet points, and sentiment scores in the article's native language. Thai text classification into seven specialized categories is executed through a hybrid pipeline combining contextual keyword-based cues with machine learning models (TF-IDF with SVM and WangchanBERTa). The frontend is built as a responsive Single Page Application (SPA) using Vanilla JavaScript ES modules and Tailwind CSS, communicating bidirectionally with the server via WebSocket (Socket.IO). It provides an in-app modal reader, trending topic visualization, and privacy-compliant personalized recommendations based on aggregate engagement metrics under PDPA standards. The evaluation results demonstrate that the pipeline operates with high stability, low latency, and significantly reduces the effort required for readers to follow diverse news outlets.